\subsection{Recollection: Bousfield localization}
\label{subsec:2dloc-bousfield}

\gutentag{0028}%
We recall the one-dimensional situation that we categorify.

\begin{defi}
\gutentag{0029}%
A \emph{Bousfield localization} of an $ \infty $-category $ C $
is a localization $ L: C \to D $ in the $ (\infty, 2) $-category of $ \infty $-categories.
\end{defi}

\gutentag{002A}%
A Bousfield localization can be described as a category obtained from $ C $ by formally
inverting a class of morphisms.

\subsubsection*{Local objects}
\gutentag{002B}%
Let $\C$ be an $\infty$-category and $S$ a class of morphisms of $\C$. An object $c$ is \emph{$S$-local} if $s^{\natural}\colon \mapp_{\C}(y,c) \to \mapp_{\C}(x,c)$ is an equivalence for every $(s\colon x \to y) \in S$; let $\iota\colon \C_{S} \subset \C$ be the full subcategory of $S$-local objects.
By the Yoneda lemma, $s$ gives a functor $R_{s}\colon \C \to \func(\Delta^{1},\Spc)$,
sending $ c $ to $s^{\natural}\colon \mapp_{\C}(y,c) \Rightarrow \mapp_{\C}(x,c)$.
The inclusion of the equivalences $\Spc \subset \func(\Delta^{1},\Spc)$ is fully faithful, and pullbacks of fully faithful functors are fully faithful, so
\begin{equation}
\label{eq:bousfield-pullback}
\C_{S} \simeq \C \times_{\prod_{S}\func(\Delta^{1},\Spc)} \textstyle\prod_{S}\Spc \ .
\end{equation}
Its two-dimensional form will be \cref{thm:local-pullback}. Here the walking equivalence is contractible, so being an equivalence is a property and there is no condition on morphisms.

\subsubsection*{The presentable case}
\gutentag{002C}%
Let $\Edge_{1} := \func(\Delta^{1},\Spc) \simeq \Free(\Delta^{1})$ (presheaves of spaces on $\Delta^{1} \simeq (\Delta^{1})\opcat$) be the free presentable $\infty$-category on an arrow, so $\func^{L}(\Edge_{1},\C) \simeq \func(\Delta^{1},\C)$ for $\C \in \prl$. The free presentable $\infty$-category on an \emph{equivalence} is $\Eqc := \Spc$.
The functor $\Delta^{1} \to \mathrm{pt}$ induces $e_{!}\colon \Edge_{1} \to \Eqc$, $(A \to B) \mapsto B$, with fully faithful right adjoint $e^{*}\colon B \mapsto (B = B)$.

\gutentag{002D}%
\reviewcomment{I think it's better to put it this way: if the left adjoint exists, then inverting $ S $ is a Bousfield localization;
in the presentable case, we can guarentee the existence of the left adjoint if $ S $ is strongly saturated and small generated;
and in this case the pushout in catinf turns out to be presentable, so there is no ambiguity what pushout we take.}
\begin{prop}
\label{prop:bousfield-pushout}
\gutentag{002E}%
Let $\C \in \prl$ and let $S$ be a set of morphisms of $\C$, viewed as a map $S\colon \bigoplus_{S}\Edge_{1} \to \C$ in $\prl$ ($\bigoplus$ the coproduct in $\prl$).
\begin{enumerate}
\item The pushout
\[
\C[S^{-1}] := \C \sqcup_{\bigoplus_{S}\Edge_{1}} \textstyle\bigoplus_{S}\Eqc
\]
in $\prl$ satisfies
\[
\func^{L}(\C[S^{-1}],\D) \simeq \{F \in \func^{L}(\C,\D) : F(s) \text{ invertible for all } s \in S\} \ .
\]
\item The right adjoint of $L\colon \C \to \C[S^{-1}]$ is the inclusion $\iota\colon \C_{S} \subset \C$. So $\C_{S} \simeq \C[S^{-1}]$ is reflective.
\item $T := \iota L$ is an idempotent monad on $\C$ and $\C_{S} \simeq \Alg_{T}(\C)$.
\end{enumerate}
\end{prop}

\begin{proof}
\gutentag{002F}%
(1) holds since $\func^{L}(\Edge_{1},\D) \simeq \func(\Delta^{1},\D)$ and $\func^{L}(\Eqc,\D) \simeq \D$ is the full subcategory of equivalences.
(2) Under $\prl \simeq (\prr)\opcat$ \cite[5.5.3.4]{htt} the pushout becomes a pullback in $\prr$, computed in $\widehat{\Cat}$ \cite[5.5.3.18]{htt}. The right adjoint of $s\colon \Edge_{1} \to \C$ is $R_{s}$ and that of $e_{!}$ is $e^{*}$, so this pullback is \eqref{eq:bousfield-pullback}.
(3) As $\iota$ is fully faithful, the counit $L\iota \Rightarrow \mathrm{id}$ and hence $\mu = \iota\varepsilon L$ are equivalences; a fully faithful right adjoint is monadic.
\end{proof}

\gutentag{0030}%
Only the existence of $L$ uses presentability of $\C$ and smallness of $S$.
The categorification will follow the same split. It will replace the walking equivalence (contractible) by the walking reflection and coreflection of \cref{subsec:mates-morphisms} (not contractible, but being reflective or coreflective is still a property by \cref{lem:mates-refl}). Local objects will be defined and studied in an arbitrary $(\infty,2)$-category in \cref{subsec:2dloc-local,subsec:2dloc-plain}, saturated classes in \cref{subsec:2dloc-saturation}, and presentability will be used for the existence of the left adjoint in \cref{subsec:2dloc-presentable}.
The local objects will form only a \emph{locally full} subcategory, and the monad will not be idempotent.
